%----------------------------------------------------------------------------------------
%  Kangana Ajay Khurana - Resume (HolboxAI :: AI Research Engineer / Scientist)
%
%  HOW TO USE ON OVERLEAF
%   1. overleaf.com -> New Project -> Blank Project
%   2. Replace the contents of main.tex with this file
%   3. Menu -> Compiler: pdfLaTeX -> Recompile
%
%  Everything inside [square brackets] is a placeholder to fill in before sending.
%
%  HONESTY NOTE (delete this comment before submitting):
%  Every claim below either comes from your own GitHub repository
%  (some1iu/movie-streaming-platform, inspected line by line), from the team repository you
%  linked (chetisha28/incursion-acm), or from information you supplied about your own work.
%  No PyTorch training, Transformer implementation, RLHF/DPO/GRPO work, or published research
%  is claimed, because no evidence for it exists yet. Once you have code, add it as a project
%  with a repository link and move the corresponding keyword out of "Research Interests".
%----------------------------------------------------------------------------------------
\documentclass[10pt,a4paper]{article}

\usepackage[T1]{fontenc}
\usepackage[utf8]{inputenc}
\usepackage{lmodern}
\usepackage[a4paper,top=1.15cm,bottom=1.0cm,left=1.5cm,right=1.5cm]{geometry}
\usepackage{enumitem}
\usepackage{xcolor}
\usepackage{microtype}
\usepackage{hyperref}

\definecolor{accent}{HTML}{1F3A5F}
\definecolor{rulegray}{HTML}{9AA5B1}
\hypersetup{colorlinks=true,urlcolor=accent,linkcolor=accent}

\pagestyle{empty}
\setlength{\parindent}{0pt}
\setlength{\parskip}{2pt}
\setlist[itemize]{leftmargin=1.1em,itemsep=1.5pt,topsep=2pt,parsep=0pt,label=\textbullet}

% Body text is set one step smaller than 10pt so the page stays clean and single-sided.
\AtBeginDocument{\small}

% Section heading: small caps + hairline rule
\newcommand{\cvsection}[1]{%
  \par\vspace{7pt}%
  {\fontsize{11}{13}\selectfont\bfseries\scshape\color{accent}#1\par}%
  \vspace{1.5pt}%
  {\color{rulegray}\rule{\textwidth}{0.55pt}\par}%
  \vspace{2.5pt}%
}

% Entry heading: title on the left, right-aligned meta on the right
\newcommand{\entry}[2]{%
  \par{\textbf{#1}}\hfill{\itshape\color{black!65}#2}\par}

% Grey sub-line under an entry
\newcommand{\subtle}[1]{%
  \par{\color{black!65}#1}\par\vspace{1pt}}

\begin{document}

%--------------------------------------- HEADER ----------------------------------------
\begin{center}
  {\fontsize{20}{22}\selectfont\bfseries\scshape\color{accent} Kangana Ajay Khurana}\\[2pt]
  {\small Computer Science \& Engineering Undergraduate \;\textbullet\; [City], India}\\[1pt]
  {\small
    \href{mailto:kanganakhurana04@gmail.com}{kanganakhurana04@gmail.com} \;\textbar\;
    \href{https://github.com/some1iu}{github.com/some1iu} \;\textbar\;
    \href{https://www.linkedin.com/in/PLACEHOLDER}{linkedin.com/in/PLACEHOLDER} \;\textbar\;
    [+91 XX XXX XXXXX]
  }
\end{center}

%--------------------------------------- SUMMARY ----------------------------------------
\cvsection{Summary}
Computer Science undergraduate aiming at research-oriented AI engineering. Foundation so far: a
self-written Python and SQL application, a four-person full-stack web project, and coursework in
algorithms, discrete mathematics, probability, object-oriented programming and databases. Currently
strengthening linear algebra, probability and optimisation in order to implement learning algorithms
from first principles, with a strong interest in deep learning, reinforcement learning and LLM
post-training.

%------------------------------------- INTERESTS ----------------------------------------
\cvsection{Research Interests}
Deep learning optimisation \;\textbullet\; reinforcement learning and post-training (reward modelling,
RLHF, DPO, GRPO) \;\textbullet\; LLM and VLM architectures, multimodal learning \;\textbullet\;
efficiency topics: Mixture-of-Experts, state-space models. \emph{(Study interests; not yet hands-on.)}

%------------------------------------- EDUCATION ----------------------------------------
\cvsection{Education}
\entry{[B.Tech / B.E.] in Computer Science and Engineering --- [University Name]}
      {[Start Year] -- [End Year]}
\subtle{[City], India \;\textbullet\; Current year: [e.g.\ First Year] \;\textbullet\;
        Expected graduation: [Month Year] \;\textbullet\; [CGPA / percentage, optional]}

\vspace{2pt}
\entry{Senior Secondary (Class XII), CBSE --- Delhi Public School, Pune}{2024 -- 2025}
\subtle{Computer Science project track; the Python project below was written for this course}

\vspace{2pt}
\textbf{Coursework and current study:} Data Structures \& Algorithms; Discrete Mathematics and
Probability; Object-Oriented Programming; Database Management Systems and relational SQL;
[add others, e.g.\ Linear Algebra, Computer Organisation].

%--------------------------------------- SKILLS ----------------------------------------
\cvsection{Technical Skills}
\begin{itemize}
  \item \textbf{Programming:} Python (project experience); C and C++ (coursework); SQL (joins,
        subqueries, aggregate functions, grouping, filtering).
  \item \textbf{Web and backend:} HTML and CSS (fundamentals); JavaScript (basics); Node.js and
        Express (recently started); React/Next.js client work in a team project; consuming and
        debugging REST APIs.
  \item \textbf{Databases and tools:} MySQL; PostgreSQL with Prisma (schema, migrations, seeding);
        Git and GitHub; VS Code.
  \item \textbf{Currently learning:} NumPy (in progress) and introductory Pandas; Node.js backend
        development; linear algebra, probability and optimisation problem sets.
\end{itemize}

%--------------------------------------- PROJECTS ----------------------------------------
\cvsection{Projects}

\entry{Movie Streaming Platform}{Python, MySQL}
\subtle{Individual project \;\textbullet\; console application, single file, $\approx$820 lines
        \;\textbullet\;
        \href{https://github.com/some1iu/movie-streaming-platform}%
             {github.com/some1iu/movie-streaming-platform}}
\begin{itemize}
  \item Wrote a menu-driven console application on top of MySQL Connector/Python that handles account
        creation and login, per-user genre preferences, movie search by title/genre/year, and rent or
        purchase transactions.
  \item Built a filter-driven SQL query builder that appends \texttt{WHERE} clauses only for the
        criteria the user supplies, using parameterised queries for the search path; modelled the data
        in three tables (user, movie, bill) and committed every transaction as a bill record.
  \item Implemented genre-based recommendations, retrieval of binary movie data from the database and
        itemised bill generation through file handling; centralised error handling for input, type, key
        and \texttt{mysql.connector} errors so the session exits gracefully instead of crashing.
\end{itemize}

\vspace{3pt}
\entry{ACM Incursion --- Conflict-Aware Events \& Expenses Manager}{Team of four}
\subtle{Next.js/React client, Express API, Prisma, PostgreSQL \;\textbullet\;
        \href{https://github.com/chetisha28/incursion-acm}%
             {github.com/chetisha28/incursion-acm}}
\begin{itemize}
  \item Contributed to a full-stack platform for ACM student chapters that flags venue, time and
        audience collisions before an event is published, gates proposals behind peer endorsement, and
        tracks expenses against semester budgets.
  \item Worked on the Next.js App Router client and its REST integration, and contributed to parts of
        the Express/Prisma backend: JWT-protected endpoints, role-based access across viewer, executive
        and core tiers, Zod request validation and \texttt{Decimal(12,2)} money columns.
  \item Debugged integration problems across Prisma migrations, seed data and CORS configuration to get
        the client running against the API locally.
\end{itemize}

\vspace{3pt}
\entry{Open-Source Codebase Exploration --- Typebot}{Bun, Nx, Prisma, PostgreSQL}
\subtle{Exploration only, no upstream contribution \;\textbullet\;
        \href{https://github.com/some1iu/typebot.io}{github.com/some1iu/typebot.io}}
\begin{itemize}
  \item Independently set up and explored a large self-hosted chatbot-builder monorepo: Bun and Nx
        workspace configuration, package and dependency graph, Prisma client generation and migrations,
        environment variables, and the PostgreSQL services required for local development; worked
        through dev-environment failures until the stack ran.
\end{itemize}

%------------------------------------- FOUNDATIONS ----------------------------------------
\cvsection{Academic Foundations}
\begin{itemize}
  \item \textbf{Algorithms:} backtracking and recursion, working systematically through Sartaj Sahni's
        algorithms text; standard data-structure and complexity exercises.
  \item \textbf{Mathematics:} discrete mathematics and probability problem sets; coding-theory problems
        including binary symmetric channels and transmission probabilities.
  \item \textbf{Object-oriented programming:} compile-time polymorphism (overloading) versus runtime
        polymorphism (overriding); class design and inheritance exercises.
  \item \textbf{Databases:} relational schema design and SQL practice --- joins, subqueries, grouping,
        aggregates and filtering.
\end{itemize}

\end{document}
